% !TEX root =  main.tex

\chapter{Introduction}
\label{chap:introduction}
\pagestyle{plain}
\setcounter{page}{1}

\section{Problem Statement}

Query optimization in relational databases hinges on cardinality estimation---predicting how many rows a query will return before actually executing it. Get the estimate right and the optimizer picks a fast execution plan. Get it wrong and the plan can be orders of magnitude slower than necessary. Traditional estimators lean on assumptions like attribute independence and uniform value distributions that rarely hold in practice, and the resulting errors compound across joins, leading to plans that are far from optimal.

Over the past few years, learned cardinality estimators have shown that neural networks can do a better job. The Multi-Set Convolutional Network (MSCN) by Kipf et al.~\cite{DBLP:conf/cidr/KipfKRLBK19} demonstrated that a relatively simple deep learning model can capture cross-table correlations that rule-based estimators miss entirely. But learned models come with their own cost: they need large quantities of labeled training data, and each label requires executing a query on a live database to record its true row count. On a complex schema this execution-based labeling can take hours, making it the dominant bottleneck in any pipeline that trains learned estimators from scratch.

Meanwhile, large language models have become remarkably good at generating SQL. Work like SQLStorm~\cite{DBLP:journals/pvldb/SchmidtLBN25} has shown that LLMs can produce diverse, realistic query workloads at scale when given a schema and a well-crafted prompt. But generating queries is only half the problem. Those queries still need to be executed to obtain labels, and that execution cost remains.

The question this thesis addresses is: can we close the loop---generate a large, diverse pool of queries with an LLM, then use active learning to pick only the most informative ones for labeling---so that we get a good estimator without paying the full cost of executing every query in the pool?

\section{Proposed Approach}

We treat the LLM not as a runtime component but as a data factory. Given a database schema, the LLM generates thousands of synthetic \texttt{SELECT COUNT(*)} queries. These queries are validated, filtered, and encoded into the tensor format that MSCN expects. Then, instead of labeling them all, we use active learning to identify the queries that the model is most uncertain about and label only those. The model trains on the labeled subset, its uncertainty shifts, and the next round of selection targets different queries. After several rounds the model has seen a compact but highly informative training set.

The whole process is packaged as an automated, schema-agnostic pipeline. A researcher points it at a database, specifies a few hyperparameters, and walks away. The pipeline generates queries, validates them through four filtering layers, labels the selected subset, trains the MSCN model, and outputs a full evaluation dashboard. No manual query writing, no template engineering, no human in the loop.

The motivation draws on two bodies of prior work. Wang et al.~\cite{DBLP:journals/pvldb/WangQWWZ21} showed that the diversity and structural properties of the training workload have a profound effect on learned estimator quality, and Davitkova and Michel~\cite{DBLP:journals/corr/abs-2512-20271} demonstrated that generative AI can automate the creation of training data for database components. We build on both: using LLMs for diverse workload generation and active learning for efficient labeling.

\section{Contributions}

This thesis makes four contributions. First, it demonstrates that LLMs can serve as effective workload generators for learned cardinality estimation, producing queries with structural diversity that rivals or exceeds what template-based methods achieve. Second, it integrates active learning into the training loop, showing that uncertainty-based query selection reduces the labeling budget needed to reach a given accuracy level. Third, it packages everything into a fully automated pipeline that works out of the box on any schema with defined foreign keys, lowering the barrier to entry for researchers who want to experiment with learned estimators. Fourth, it provides an extensive experimental evaluation on the IMDB benchmark that quantifies the trade-offs between different acquisition strategies, workload sizes, and training configurations.